
\documentclass[10pt,a4paper]{article}
\usepackage[utf8]{inputenc}
\usepackage[french]{babel}
\usepackage[T1]{fontenc}            
\usepackage{lmodern}
\usepackage[top=2.5cm, bottom=2.5cm, left=2.5cm, right=2.5cm]{geometry}
\usepackage{amsmath}
\usepackage{amsfonts}
\usepackage{amssymb}
\usepackage{amsthm}
\usepackage{calrsfs}
\usepackage{stmaryrd}
\usepackage{minted}
\usepackage{fancyhdr}
\usepackage[inline]{asymptote} 
\usepackage[linesnumbered,ruled,french,onelanguage]{algorithm2e}
\SetKwComment{Comment}{$\triangleright$\ }{}
\usepackage{hyperref}
\hypersetup{
    colorlinks,
    linktoc=section,
    linkcolor=red,
    urlcolor=orange,
    filecolor=red
}
\usepackage{times}
\usepackage{todonotes}
\usepackage{latexsym}
\usepackage{verbatim}
\usepackage{amsmath,amssymb}
\usepackage{upgreek}
\usepackage{dsfont}
\usepackage{framed}
\usepackage{amsfonts}
\newcommand{\tocheck}[1]{\textcolor{blue}{#1}}
\newcommand{\overbar}[1]{\mkern 3mu\overline{\mkern-3mu#1\mkern-3mu}\mkern 3mu}
\usepackage{fullpage,setspace}
\linespread{1.5}
\usepackage{enumerate}
\usepackage{mathrsfs}
\usepackage{amsthm}
\theoremstyle{definition}
\newtheorem{lemma}{Lemme}
\newtheorem{theorem}{Théorème}
\newtheorem{observation}[theorem]{Observation}
\newtheorem{definition}{Définition}
\newtheorem{heuristique}{Heuristique}
\newtheorem{scenario}{Scénario}
\newtheorem{fact}{Fait}
\newtheorem{example}{Exemple}
\newtheorem{postulate}{Postulat}
\newtheorem{proposition}{Proposition}
\newtheorem{corollary}{Corollaire}
\newtheorem{property}{Propriété}
\numberwithin{lemma}{subsection}
\numberwithin{theorem}{subsection}
\numberwithin{definition}{subsection}
\numberwithin{proposition}{subsection}
\numberwithin{corollary}{subsection}
\numberwithin{property}{subsection}
\numberwithin{example}{subsection}
\numberwithin{heuristique}{subsection}
\numberwithin{scenario}{subsection}
\newenvironment{proofi} {\noindent\emph{Proof:}} {\hfill $\square$\vspace{0.2cm}}
\newcommand{\eod}{\ensuremath{\hfill\dashv}}
\newcommand{\eop}{\ensuremath{\hfill\clubsuit}}
\newcommand{\fun}[1]{\ensuremath{\mbox{\it #1}}}
\newcommand{\expo}[2]{{#1}^{\mbox{\scriptsize \sl #2}}}
\newcommand{\var}{\ensuremath{\mathbf{var}}\xspace}
\newcommand{\exv}{\ensuremath{\mathbf{exv}}\xspace}
\newcommand{\dsv}{\ensuremath{\mathbf{dsv}}\xspace}
\newcommand{\term}{\ensuremath{\mathbf{term}}\xspace}
\newcommand{\pred}{\ensuremath{\mathbf{pred}}\xspace}
\newcommand{\fr}{\ensuremath{\mathbf{fr}}\xspace}
\newcommand{\cnst}{\ensuremath{\mathbf{cnst}}\xspace}
\newcommand{\Vars}{\ensuremath{\mathsf{Vars}}\xspace}
\newcommand{\Pvars}{\ensuremath{\mathsf{Pvars}}\xspace}
\newcommand{\Fvars}{\ensuremath{\mathsf{Fvars}}\xspace}
\newcommand{\Terms}{\ensuremath{\mathsf{Terms}}\xspace}
\newcommand{\Cons}{\ensuremath{\mathsf{Const}}\xspace}
\newcommand{\Preds}{\ensuremath{\mathsf{Preds}}\xspace}
\newcommand{\ic}{\ensuremath{\cdot\{}\xspace}
\newcommand{\ci}{\ensuremath{\}\cdot}\xspace}
\newcommand{\id}{\ensuremath{:\hspace{-1.56mm}\{}\xspace}
\newcommand{\di}{\ensuremath{\}\hspace{-1.56mm}:}\xspace}
\newcommand{\bet}{\ensuremath{\beta^{\diamond}}\xspace}
\newcommand{\sep}{\ensuremath{\mathbf{sep}}\xspace}
\newcommand{\vect}[1]{\mathbf{#1}}
\newcommand{\vars}[1]{\fun{vars}{(#1)}}
\newcommand{\terms}[1]{\fun{terms}{(#1)}}
\newcommand{\atoms}[1]{\fun{atoms}{(#1)}}
\newcommand{\type}[1]{\fun{type}{(#1)}}
\newcommand{\const}{a}
\newcommand{\atom}{\alpha}

\newcommand{\Zero}{\includegraphics[height=12px]{img/jetons/0}}
\newcommand{\Successor}{\includegraphics[height=12px]{img/jetons/S}}
\newcommand{\Identity}{\includegraphics[height=12px]{img/jetons/I}}
\newcommand{\Left}{\includegraphics[height=12px]{img/jetons/gauche}}
\newcommand{\Right}{\includegraphics[height=12px]{img/jetons/droite}}
\newcommand{\Composition}{\includegraphics[height=12px]{img/jetons/composition}}
\newcommand{\Recurse}{\includegraphics[height=12px]{img/jetons/R}}
\newcommand{\Mu}{\includegraphics[height=12px]{img/jetons/mu}}

\newcounter{deuxfois}
\newcommand{\sig}{\ensuremath{\mathbf{sig}}\xspace}

\usepackage{color}
\definecolor{gren}{rgb}{0.158, 0.488, 0.408}
\definecolor{bren}{rgb}{0.45, 0.21, 0.000}

\title{Notes de cours d'informatique théorique (HMIN 118)}
\author{Guillaume Pérution-Kihli, Julien Rodriguez}


\begin{document}

\maketitle
\renewcommand{\contentsname}{Sommaire}
\tableofcontents
\pagebreak

\section{Théorie de l'information de Shannon}
	\subsection{Mesure de l'information de Hartley}
	
L'information est le nombre de réponses possible.
\par Hartley suppose que pour composer un message, il faut avoir un alphabet dans lequel le composer. Cet alphabet doit comporter un nombre fini $s$ de symboles. Un message comportant $n$ symboles tirés d'un alphabet de taille $s$ n'est donc qu'un parmi les $s^n$ messages possibles.
\par D'après Hartley, la quantité d'information intrinsèque au message doit être proportionnelle à sa longueur fois une quantité qui ne dépend que du nombre de symboles de l'alphabet. Autrement dit, la quantité d'information $H$ d'un message de longueur $n$ est de la forme $H = k * n$ où la constante $k$ est à déterminer.
\par Pour obtenir la fonction désirée, Hartley propose d'utiliser le logarithme et de calculer $$H = log(s^n) = n*log(s) $$.
\par Ce logarithme peut avoir différentes bases, et l'unité de mesure de l'information change en fonction de ces bases.
\par
\begin{center}
\begin{tabular}{|c|c|c|}
\hline
base &  & nom de l'unité \\
\hline
$2$ & $2 (log_2)$ & bit \\
\hline
$e$ & $e (ln)$ & nat \\
\hline
$10$ & $10(log_{10})$ & hartley \\
\hline
\end{tabular}
\end{center}

\par Cette méthode de mesure a cependant un important défaut : tous les symboles du message sont considérés comme contenant la même quantité d'information puisqu'on les multiplie tous par la même constante.

	\subsection{Mesure de l'information de Shannon}

Pour Shannon, l'unité d'information est le bit. Avec le bit comme mesure, l’information devient une notion essentiellement probabiliste donc quantifiable. Et s’il y a beaucoup de manières de passer un message en binaire, il n’y en a qu’une qui permette de réaliser la quantité d’information de Shannon. L’entropie de Shannon mesure la quantité d’information d’un message : un signal peu informatif est redondant et donc prédictible d’un point de vue des probabilités, alors qu’un signal très informatif est très diversifié et donc peu prédictible.
\par Pour une source, qui est une variable aléatoire discrète $X$ comportant $n$ symboles, chaque symbole $x_i$ ayant une probabilité $P_i$ d'apparaître, l'entropie $H$ de la source $X$ est définie comme :
$$H_b(X)= -\mathbb E [\log_b {P(X)}] = \sum_{i=1}^nP_i\log_b \left(\frac{1}{P_i}\right)=-\sum_{i=1}^nP_i\log_b P_i.\,\!$$
où $\mathbb  E$ désigne l'espérance mathématique, et $\log _{b}$ le logarithme en base $b$. On utilise en général un logarithme à base 2 car l'entropie possède alors les unités de bit/symbole. Les symboles représentent les réalisations possibles de la variable aléatoire $X$. Dans ce cas, on peut interpréter $H(X)$ comme le nombre de questions à réponse oui/non que doit poser en moyenne le récepteur à la source, ou la quantité d'information en bits que la source doit fournir au récepteur pour que ce dernier puisse déterminer sans ambiguïté la valeur de $X$.
$$H(X)=H_2(X)= -\sum_{i=1}^nP_i\log_2 P_i.\,\!$$

\begin{example}\label{ex:entropie_shannon}
Prenons $X = yabadabadoo$. On obtient la répartition qui suit :
\par\begin{tabular}{|c|c|c|c|c|c|}
\hline
Symbole & a & b & d & o & y \\
\hline
Nombre d'occurrences & 4 & 2 & 2 & 2 & 1 \\
\hline
Probabilité d'apparition $P_i$ & $\frac{4}{11}$ & $\frac{2}{11}$ & $\frac{2}{11}$ & $\frac{2}{11}$ & $\frac{1}{11}$ \\
\hline
\end{tabular}
\par L'entropie de $X$ est :
$$H(X) = -\sum_{i=1}^5P_i\log_2 P_i = -(\frac{4}{11}\log_2\frac{4}{11}+3*
(\frac{2}{11}\log_2\frac{2}{11})+
\frac{1}{11}\log_2\frac{1}{11}) \simeq 2.19$$
\end{example}

\section{Codes et techniques de compression}
	\subsection{Codes, codes préfixes et inégalité de Kraft}

\begin{scenario}
A veut transmettre quelques informations à B.
\end{scenario}

$element \in \chi \rightarrow chaine ~binaire$

\begin{definition}{Fonction de décodage}
\par $D : {0,1}^* \rightarrow \chi$.
\par $E = D^{-1}$.
\end{definition}

\begin{theorem}{Inégalité de Kraft}
\par Il existe un code uniquement décodable sur un alphabet de taille $r$, avec $n$ mots de code de tailles $l_1, l_2, ..., l_n$, ssi $$\sum_{i=1}^n r^{-l_i} \leqslant 1$$
\end{theorem}

\begin{definition}{Sans préfixe}
\par Un code est dit sans préfixe si aucun mot de code n'est préfixe d'un autre.
\end{definition}

	\subsection{Codage de Shannon-Fano}
	
Le codage de Shannon-Fano est un algorithme de compression de données sans perte élaboré par Robert Fano à partir d'une idée de Claude Shannon.
Il s'agit d'un codage entropique produisant un code préfixe très similaire à un code de Huffman, bien que pas toujours optimal, contrairement à ce dernier.

\begin{algorithm}[H]
\caption{Algorithme de Shannon-Fano}
Pour chaque symbole, compter la fréquence d'apparition\;
Les ordonner des plus fréquents aux moins fréquents \;
Diviser la liste de symboles en deux parties de façon à avoir des totaux proches \;
Assigner 0 à la partie gauche de la division, et 1 à la partie droite \;
Recommencer les étapes 3 et 4\;
\end{algorithm}

\begin{example}
Reprenons l'exemple \ref{ex:entropie_shannon}. On obtient la liste de symboles ordonnés : $<(a:4), (b:2), (d:2), (o:2), (y:1)>$.
\par On va diviser cette liste une première fois en deux listes ayant à peu près les mêmes nombres d'occurrences, qu'on va représenter sous forme d'arbre :
%:-+-+-+- Engendré par : http://math.et.info.free.fr/TikZ/Arbre/
\begin{center}
% Racine en Haut, développement vers le bas
\begin{tikzpicture}[xscale=1,yscale=1]
% Styles (MODIFIABLES)
\tikzstyle{fleche}=[->,>=latex,thick]
\tikzstyle{noeud}=[fill=lightgray,rectangle,draw]
\tikzstyle{feuille}=[fill=lightgray,rectangle,draw]
\tikzstyle{etiquette}=[midway,fill=white]
% Dimensions (MODIFIABLES)
\def\DistanceInterNiveaux{3}
\def\DistanceInterFeuilles{2}
% Dimensions calculées (NON MODIFIABLES)
\def\NiveauA{(-0)*\DistanceInterNiveaux}
\def\NiveauB{(-0.7)*\DistanceInterNiveaux}
\def\InterFeuilles{(1)*\DistanceInterFeuilles}
% Noeuds (MODIFIABLES : Styles et Coefficients d'InterFeuilles)
\node[noeud] (R) at ({(1)*\InterFeuilles},{\NiveauA}) {$11$};
\node[feuille] (Ra) at ({(0)*\InterFeuilles},{\NiveauB}) {$6: (a:4), (b:2)$};
\node[feuille] (Rb) at ({(2)*\InterFeuilles},{\NiveauB}) {$5: (d:2), (o:2), (y:1)$};
% Arcs (MODIFIABLES : Styles)
\draw[fleche] (R)--(Ra) node[etiquette] {$0$};
\draw[fleche] (R)--(Rb) node[etiquette] {$1$};
\end{tikzpicture}
\end{center}
On répète le processus jusqu'à obtenir un arbre où chaque feuille représente un symbole :
%:-+-+-+- Engendré par : http://math.et.info.free.fr/TikZ/Arbre/
\begin{center}
% Racine en Haut, développement vers le bas
\begin{tikzpicture}[xscale=1,yscale=1]
% Styles (MODIFIABLES)
\tikzstyle{fleche}=[->,>=latex,thick]
\tikzstyle{noeud}=[fill=lightgray,rectangle,draw]
\tikzstyle{feuille}=[fill=lightgray,rectangle,draw]
\tikzstyle{etiquette}=[midway,fill=white]
% Dimensions (MODIFIABLES)
\def\DistanceInterNiveaux{3}
\def\DistanceInterFeuilles{2}
% Dimensions calculées (NON MODIFIABLES)
\def\NiveauA{(-0)*\DistanceInterNiveaux}
\def\NiveauB{(-0.7)*\DistanceInterNiveaux}
\def\NiveauC{(-1.4)*\DistanceInterNiveaux}
\def\InterFeuilles{(1)*\DistanceInterFeuilles}
% Noeuds (MODIFIABLES : Styles et Coefficients d'InterFeuilles)
\node[noeud] (R) at ({(1.5)*\InterFeuilles},{\NiveauA}) {$11$};
\node[noeud] (Ra) at ({(0.5)*\InterFeuilles},{\NiveauB}) {$6$};
\node[feuille] (Raa) at ({(0)*\InterFeuilles},{\NiveauC}) {$a$};
\node[feuille] (Rab) at ({(1)*\InterFeuilles},{\NiveauC}) {$b$};
\node[noeud] (Rb) at ({(2.5)*\InterFeuilles},{\NiveauB}) {$5$};
\node[feuille] (Rba) at ({(2)*\InterFeuilles},{\NiveauC}) {$3: (d:2), (y:1)$};
\node[feuille] (Rbb) at ({(3)*\InterFeuilles},{\NiveauC}) {$o$};
% Arcs (MODIFIABLES : Styles)
\draw[fleche] (R)--(Ra) node[etiquette] {$0$};
\draw[fleche] (Ra)--(Raa) node[etiquette] {$0$};
\draw[fleche] (Ra)--(Rab) node[etiquette] {$1$};
\draw[fleche] (R)--(Rb) node[etiquette] {$1$};
\draw[fleche] (Rb)--(Rba) node[etiquette] {$0$};
\draw[fleche] (Rb)--(Rbb) node[etiquette] {$1$};
\end{tikzpicture}
\end{center}
%:-+-+-+- Engendré par : http://math.et.info.free.fr/TikZ/Arbre/
\begin{center}
% Racine en Haut, développement vers le bas
\begin{tikzpicture}[xscale=1,yscale=1]
% Styles (MODIFIABLES)
\tikzstyle{fleche}=[->,>=latex,thick]
\tikzstyle{noeud}=[fill=lightgray,rectangle,draw]
\tikzstyle{feuille}=[fill=lightgray,rectangle,draw]
\tikzstyle{etiquette}=[midway,fill=white]
% Dimensions (MODIFIABLES)
\def\DistanceInterNiveaux{3}
\def\DistanceInterFeuilles{2}
% Dimensions calculées (NON MODIFIABLES)
\def\NiveauA{(-0)*\DistanceInterNiveaux}
\def\NiveauB{(-0.7)*\DistanceInterNiveaux}
\def\NiveauC{(-1.4)*\DistanceInterNiveaux}
\def\NiveauD{(-2.1)*\DistanceInterNiveaux}
\def\InterFeuilles{(1)*\DistanceInterFeuilles}
% Noeuds (MODIFIABLES : Styles et Coefficients d'InterFeuilles)
\node[noeud] (R) at ({(2)*\InterFeuilles},{\NiveauA}) {$11$};
\node[noeud] (Ra) at ({(0.5)*\InterFeuilles},{\NiveauB}) {$6$};
\node[feuille] (Raa) at ({(0)*\InterFeuilles},{\NiveauC}) {$a$};
\node[feuille] (Rab) at ({(1)*\InterFeuilles},{\NiveauC}) {$b$};
\node[noeud] (Rb) at ({(3)*\InterFeuilles},{\NiveauB}) {$5$};
\node[noeud] (Rba) at ({(2.5)*\InterFeuilles},{\NiveauC}) {$3$};
\node[feuille] (Rbaa) at ({(2)*\InterFeuilles},{\NiveauD}) {$d$};
\node[feuille] (Rbab) at ({(3)*\InterFeuilles},{\NiveauD}) {$y$};
\node[feuille] (Rbb) at ({(4)*\InterFeuilles},{\NiveauC}) {$o$};
% Arcs (MODIFIABLES : Styles)
\draw[fleche] (R)--(Ra) node[etiquette] {$0$};
\draw[fleche] (Ra)--(Raa) node[etiquette] {$0$};
\draw[fleche] (Ra)--(Rab) node[etiquette] {$1$};
\draw[fleche] (R)--(Rb) node[etiquette] {$1$};
\draw[fleche] (Rb)--(Rba) node[etiquette] {$0$};
\draw[fleche] (Rba)--(Rbaa) node[etiquette] {$0$};
\draw[fleche] (Rba)--(Rbab) node[etiquette] {$1$};
\draw[fleche] (Rb)--(Rbb) node[etiquette] {$1$};
\end{tikzpicture}
\end{center}
On obtient ainsi des codes sans préfixe pour chaque symbole.
\par\begin{tabular}{|c|c|c|c|c|c|}
\hline
Symbole & a & b & d & o & y \\
\hline
Code & 00 & 01 & 100 & 11 & 101 \\
\hline
Taille code & 2 & 2 & 3 & 2 & 3 \\
\hline
\end{tabular}
\par La longueur totale de l'encodage de la chaîne sera : $2*4 + 2*2 + 3*2 + 2*2 + 2*1 = 25$ bits. Le nombre de symboles moyen encodés par bit est de $\frac{25}{11} \simeq 2,27 > 2,19 \simeq H(X)$.
\end{example}

	\subsection{Codage de Huffman}

L'approche du codage de Shannon-Fano est descendante : l'algorithme part de l'ensemble des symboles et divise cet ensemble récursivement jusqu'à arriver à des parties ne contenant qu'un seul symbole. L'inconvénient de cette approche est que, lorsqu'il n'est pas possible de séparer un ensemble de symboles et deux sous-ensembles de probabilités à peu près égales (c'est-à-dire lorsque l'un des sous-ensembles est beaucoup plus probable que l'autre), les codes produits ne sont pas optimaux.
\par Le codage de Huffman a une approche ascendante : l'algorithme part des symboles et regroupe ceux ayant la probabilité la plus faible, jusqu'à avoir regroupé tous les symboles. Cette approche permet d'obtenir systématiquement un code optimal au niveau du symbole, dans le pire cas de la même longueur que le code de Shannon-Fano équivalent, dans tous les autres cas plus court.

\begin{algorithm}[H]
\caption{Algorithme de Huffman}
Pour chaque symbole, créer une feuille, l'associer avec sa fréquence d'apparition et l'ajouter à une file de priorité\;
\While{il y a plus d'un nœud associé à une fréquence dans la file}
{
	Supprimer de la file les 2 nœuds ayant la fréquence la plus faible\;
	Créer un nouveau nœud ayant pour fils les 2 nœuds supprimés, l'associer à la somme des fréquences des nœuds fils\;
	Ajouter le nouveau nœud à la file\;
}
Le nœud restant dans la file est la racine de l'arbre\;
\end{algorithm}

\begin{example}
\par Reprenons l'exemple \ref{ex:entropie_shannon} avec $X = yabadabadoo$. Les étapes de construction de l'arbre de Huffman sont les suivantes :
\begin{itemize}
\item[1]
%:-+-+-+- Engendré par : http://math.et.info.free.fr/TikZ/Arbre/
\begin{center}
% Racine en Haut, développement vers le bas
\begin{tikzpicture}[xscale=1,yscale=1]
% Styles (MODIFIABLES)
\tikzstyle{fleche}=[->,>=latex,thick]
\tikzstyle{noeudTransparent}=[]
\tikzstyle{noeud}=[fill=lightgray,rectangle,draw]
\tikzstyle{feuille}=[fill=lightgray,rectangle,draw]
\tikzstyle{etiquette}=[]
% Dimensions (MODIFIABLES)
\def\DistanceInterNiveaux{3}
\def\DistanceInterFeuilles{2}
% Dimensions calculées (NON MODIFIABLES)
\def\NiveauA{(-0)*\DistanceInterNiveaux}
\def\NiveauB{(-0)*\DistanceInterNiveaux}
\def\NiveauC{(-0.5)*\DistanceInterNiveaux}
\def\InterFeuilles{(1)*\DistanceInterFeuilles}
% Noeuds (MODIFIABLES : Styles et Coefficients d'InterFeuilles)
\node[noeudTransparent] (Ra) at ({(0)*\InterFeuilles},{\NiveauB}) {$ $};
\node[feuille] (Raa) at ({(0)*\InterFeuilles},{\NiveauC}) {$a:4$};
\node[noeudTransparent] (Rb) at ({(1)*\InterFeuilles},{\NiveauB}) {$ $};
\node[feuille] (Rba) at ({(1)*\InterFeuilles},{\NiveauC}) {$b:2$};
\node[noeudTransparent] (Rc) at ({(2)*\InterFeuilles},{\NiveauB}) {$ $};
\node[feuille] (Rca) at ({(2)*\InterFeuilles},{\NiveauC}) {$d:2$};
\node[noeudTransparent] (Rd) at ({(3)*\InterFeuilles},{\NiveauB}) {$ $};
\node[feuille] (Rda) at ({(3)*\InterFeuilles},{\NiveauC}) {$o:2$};
\node[noeudTransparent] (Re) at ({(4)*\InterFeuilles},{\NiveauB}) {$ $};
\node[feuille] (Rea) at ({(4)*\InterFeuilles},{\NiveauC}) {$y:1$};
% Arcs (MODIFIABLES : Styles)
\draw[fleche] (Ra)--(Raa) node[etiquette] {$ $};
\draw[fleche] (Rb)--(Rba) node[etiquette] {$ $};
\draw[fleche] (Rc)--(Rca) node[etiquette] {$ $};
\draw[fleche] (Rd)--(Rda) node[etiquette] {$ $};
\draw[fleche] (Re)--(Rea) node[etiquette] {$ $};
\end{tikzpicture}
\end{center}
%:-+-+-+-+- Fin
\item[2]
%:-+-+-+- Engendré par : http://math.et.info.free.fr/TikZ/Arbre/
\begin{center}
% Racine en Haut, développement vers le bas
\begin{tikzpicture}[xscale=1,yscale=1]
% Styles (MODIFIABLES)
\tikzstyle{fleche}=[->,>=latex,thick]
\tikzstyle{noeudTransparent}=[]
\tikzstyle{noeud}=[fill=lightgray,rectangle,draw]
\tikzstyle{feuille}=[fill=lightgray,rectangle,draw]
\tikzstyle{etiquette}=[]
% Dimensions (MODIFIABLES)
\def\DistanceInterNiveaux{3}
\def\DistanceInterFeuilles{2}
% Dimensions calculées (NON MODIFIABLES)
\def\NiveauA{(-0)*\DistanceInterNiveaux}
\def\NiveauB{(-0)*\DistanceInterNiveaux}
\def\NiveauC{(-0.5)*\DistanceInterNiveaux}
\def\NiveauD{(-1)*\DistanceInterNiveaux}
\def\InterFeuilles{(1)*\DistanceInterFeuilles}
% Noeuds (MODIFIABLES : Styles et Coefficients d'InterFeuilles)
\node[noeudTransparent] (Ra) at ({(0)*\InterFeuilles},{\NiveauB}) {$ $};
\node[feuille] (Raa) at ({(0)*\InterFeuilles},{\NiveauC}) {$a:4$};
\node[noeudTransparent] (Rb) at ({(1)*\InterFeuilles},{\NiveauB}) {$ $};
\node[feuille] (Rba) at ({(1)*\InterFeuilles},{\NiveauC}) {$b:2$};
\node[noeudTransparent] (Rc) at ({(2)*\InterFeuilles},{\NiveauB}) {$ $};
\node[feuille] (Rca) at ({(2)*\InterFeuilles},{\NiveauC}) {$d:2$};
\node[noeudTransparent] (Rd) at ({(3.5)*\InterFeuilles},{\NiveauB}) {$ $};
\node[noeud] (Rda) at ({(3.5)*\InterFeuilles},{\NiveauC}) {$3$};
\node[feuille] (Rdaa) at ({(3)*\InterFeuilles},{\NiveauD}) {$o$};
\node[feuille] (Rdab) at ({(4)*\InterFeuilles},{\NiveauD}) {$y$};
% Arcs (MODIFIABLES : Styles)
\draw[fleche] (Ra)--(Raa) node[etiquette] {$ $};
\draw[fleche] (Rb)--(Rba) node[etiquette] {$ $};
\draw[fleche] (Rc)--(Rca) node[etiquette] {$ $};
\draw[fleche] (Rd)--(Rda) node[etiquette] {$ $};
\draw[fleche] (Rda)--(Rdaa) node[etiquette] {$ $};
\draw[fleche] (Rda)--(Rdab) node[etiquette] {$ $};
\end{tikzpicture}
\end{center}
%:-+-+-+-+- Fin
\item[3]
%:-+-+-+- Engendré par : http://math.et.info.free.fr/TikZ/Arbre/
\begin{center}
% Racine en Haut, développement vers le bas
\begin{tikzpicture}[xscale=1,yscale=1]
% Styles (MODIFIABLES)
\tikzstyle{fleche}=[->,>=latex,thick]
\tikzstyle{noeudTransparent}=[]
\tikzstyle{noeud}=[fill=lightgray,rectangle,draw]
\tikzstyle{feuille}=[fill=lightgray,rectangle,draw]
\tikzstyle{etiquette}=[]
% Dimensions (MODIFIABLES)
\def\DistanceInterNiveaux{3}
\def\DistanceInterFeuilles{2}
% Dimensions calculées (NON MODIFIABLES)
\def\NiveauA{(-0)*\DistanceInterNiveaux}
\def\NiveauB{(-0)*\DistanceInterNiveaux}
\def\NiveauC{(-0.5)*\DistanceInterNiveaux}
\def\NiveauD{(-1)*\DistanceInterNiveaux}
\def\InterFeuilles{(1)*\DistanceInterFeuilles}
% Noeuds (MODIFIABLES : Styles et Coefficients d'InterFeuilles)
\node[noeudTransparent] (Ra) at ({(0)*\InterFeuilles},{\NiveauB}) {$ $};
\node[feuille] (Raa) at ({(0)*\InterFeuilles},{\NiveauC}) {$a:4$};
\node[noeudTransparent] (Rb) at ({(1.5)*\InterFeuilles},{\NiveauB}) {$ $};
\node[noeud] (Rba) at ({(1.5)*\InterFeuilles},{\NiveauC}) {$4$};
\node[feuille] (Rbaa) at ({(1)*\InterFeuilles},{\NiveauD}) {$b$};
\node[feuille] (Rbab) at ({(2)*\InterFeuilles},{\NiveauD}) {$d$};
\node[noeudTransparent] (Rc) at ({(3.5)*\InterFeuilles},{\NiveauB}) {$ $};
\node[noeud] (Rca) at ({(3.5)*\InterFeuilles},{\NiveauC}) {$3$};
\node[feuille] (Rcaa) at ({(3)*\InterFeuilles},{\NiveauD}) {$o$};
\node[feuille] (Rcab) at ({(4)*\InterFeuilles},{\NiveauD}) {$y$};
% Arcs (MODIFIABLES : Styles)
\draw[fleche] (Ra)--(Raa) node[etiquette] {$ $};
\draw[fleche] (Rb)--(Rba) node[etiquette] {$ $};
\draw[fleche] (Rba)--(Rbaa) node[etiquette] {$ $};
\draw[fleche] (Rba)--(Rbab) node[etiquette] {$ $};
\draw[fleche] (Rc)--(Rca) node[etiquette] {$ $};
\draw[fleche] (Rca)--(Rcaa) node[etiquette] {$ $};
\draw[fleche] (Rca)--(Rcab) node[etiquette] {$ $};
\end{tikzpicture}
\end{center}
%:-+-+-+-+- Fin
\item[4]
%:-+-+-+- Engendré par : http://math.et.info.free.fr/TikZ/Arbre/
\begin{center}
% Racine en Haut, développement vers le bas
\begin{tikzpicture}[xscale=1,yscale=1]
% Styles (MODIFIABLES)
\tikzstyle{fleche}=[->,>=latex,thick]
\tikzstyle{noeudTransparent}=[]
\tikzstyle{noeud}=[fill=lightgray,rectangle,draw]
\tikzstyle{feuille}=[fill=lightgray,rectangle,draw]
\tikzstyle{etiquette}=[]
% Dimensions (MODIFIABLES)
\def\DistanceInterNiveaux{3}
\def\DistanceInterFeuilles{2}
% Dimensions calculées (NON MODIFIABLES)
\def\NiveauA{(-0)*\DistanceInterNiveaux}
\def\NiveauB{(-0)*\DistanceInterNiveaux}
\def\NiveauC{(-0.5)*\DistanceInterNiveaux}
\def\NiveauD{(-1)*\DistanceInterNiveaux}
\def\NiveauE{(-1.5)*\DistanceInterNiveaux}
\def\InterFeuilles{(1)*\DistanceInterFeuilles}
% Noeuds (MODIFIABLES : Styles et Coefficients d'InterFeuilles)
\node[noeudTransparent] (Ra) at ({(0)*\InterFeuilles},{\NiveauB}) {$ $};
\node[feuille] (Raa) at ({(0)*\InterFeuilles},{\NiveauC}) {$a:4$};
\node[noeudTransparent] (Rb) at ({(2.5)*\InterFeuilles},{\NiveauB}) {$ $};
\node[noeud] (Rba) at ({(2.5)*\InterFeuilles},{\NiveauC}) {$7$};
\node[noeud] (Rbaa) at ({(1.5)*\InterFeuilles},{\NiveauD}) {$4$};
\node[feuille] (Rbaaa) at ({(1)*\InterFeuilles},{\NiveauE}) {$b$};
\node[feuille] (Rbaab) at ({(2)*\InterFeuilles},{\NiveauE}) {$d$};
\node[noeud] (Rbab) at ({(3.5)*\InterFeuilles},{\NiveauD}) {$3$};
\node[feuille] (Rbaba) at ({(3)*\InterFeuilles},{\NiveauE}) {$o$};
\node[feuille] (Rbabb) at ({(4)*\InterFeuilles},{\NiveauE}) {$y$};
% Arcs (MODIFIABLES : Styles)
\draw[fleche] (Ra)--(Raa) node[etiquette] {$ $};
\draw[fleche] (Rb)--(Rba) node[etiquette] {$ $};
\draw[fleche] (Rba)--(Rbaa) node[etiquette] {$ $};
\draw[fleche] (Rbaa)--(Rbaaa) node[etiquette] {$ $};
\draw[fleche] (Rbaa)--(Rbaab) node[etiquette] {$ $};
\draw[fleche] (Rba)--(Rbab) node[etiquette] {$ $};
\draw[fleche] (Rbab)--(Rbaba) node[etiquette] {$ $};
\draw[fleche] (Rbab)--(Rbabb) node[etiquette] {$ $};
\end{tikzpicture}
\end{center}
%:-+-+-+-+- Fin
\item[5]%:-+-+-+- Engendré par : http://math.et.info.free.fr/TikZ/Arbre/
\begin{center}
% Racine en Haut, développement vers le bas
\begin{tikzpicture}[xscale=1,yscale=1]
% Styles (MODIFIABLES)
\tikzstyle{fleche}=[->,>=latex,thick]
\tikzstyle{etiquette}=[midway,fill=white]
\tikzstyle{noeud}=[fill=lightgray,rectangle,draw]
\tikzstyle{feuille}=[fill=lightgray,rectangle,draw]
% Dimensions (MODIFIABLES)
\def\DistanceInterNiveaux{3}
\def\DistanceInterFeuilles{2}
% Dimensions calculées (NON MODIFIABLES)
\def\NiveauA{(-0)*\DistanceInterNiveaux}
\def\NiveauB{(-0.5)*\DistanceInterNiveaux}
\def\NiveauC{(-1)*\DistanceInterNiveaux}
\def\NiveauD{(-1.5)*\DistanceInterNiveaux}
\def\InterFeuilles{(1)*\DistanceInterFeuilles}
% Noeuds (MODIFIABLES : Styles et Coefficients d'InterFeuilles)
\node[noeud] (R) at ({(1)*\InterFeuilles},{\NiveauA}) {$11$};
\node[feuille] (Ra) at ({(0)*\InterFeuilles},{\NiveauB}) {$a$};
\node[noeud] (Rb) at ({(2.5)*\InterFeuilles},{\NiveauB}) {$7$};
\node[noeud] (Rba) at ({(1.5)*\InterFeuilles},{\NiveauC}) {$4$};
\node[feuille] (Rbaa) at ({(1)*\InterFeuilles},{\NiveauD}) {$b$};
\node[feuille] (Rbab) at ({(2)*\InterFeuilles},{\NiveauD}) {$d$};
\node[noeud] (Rbb) at ({(3.5)*\InterFeuilles},{\NiveauC}) {$3$};
\node[feuille] (Rbba) at ({(3)*\InterFeuilles},{\NiveauD}) {$o$};
\node[feuille] (Rbbb) at ({(4)*\InterFeuilles},{\NiveauD}) {$y$};
% Arcs (MODIFIABLES : Styles)
\draw[fleche] (R)--(Ra) node[etiquette] {$0$};
\draw[fleche] (R)--(Rb) node[etiquette] {$1$};
\draw[fleche] (Rb)--(Rba) node[etiquette] {$0$};
\draw[fleche] (Rba)--(Rbaa) node[etiquette] {$0$};
\draw[fleche] (Rba)--(Rbab) node[etiquette] {$1$};
\draw[fleche] (Rb)--(Rbb) node[etiquette] {$1$};
\draw[fleche] (Rbb)--(Rbba) node[etiquette] {$0$};
\draw[fleche] (Rbb)--(Rbbb) node[etiquette] {$1$};
\end{tikzpicture}
\end{center}
%:-+-+-+-+- Fin
\end{itemize}
On obtient ainsi des codes sans préfixe pour chaque symbole.
\par\begin{tabular}{|c|c|c|c|c|c|}
\hline
Symbole & a & b & d & o & y \\
\hline
Code & 0 & 100 & 101 & 110 & 111 \\
\hline
Taille code & 1 & 3 & 3 & 3 & 3 \\
\hline
\end{tabular}
\par La longueur totale de l'encodage de la chaîne sera : $1*4 + 3*7 = 25$ bits. Le nombre de symboles moyen encodés par bit est de $\frac{25}{11} \simeq 2,27 > 2,19 \simeq H(X)$. On constate qu'on est sur un cas avec la même moyenne que Shannon-Fano.
\end{example}


        \par L'algorithme MTF (pour move-to-front : "déplacer vers l'avant") est un système de transformation de flot utilisé notamment dans le domaine de la compression de données en informatique. 

\subsubsection{Fonctionnement}
\par Cet algorithme consiste à remplacer chaque caractère par un indice donné par un tableau évoluant de manière dynamique. 
Le tableau est tout d'abord initialisé en rangeant les caractères utilisé pour le codage comme ceci : \\

\begin{tabular}{|c|c|c|c|c|c|c|c|c|c|c|c|c|c|c|c|}
    \hline
    Indice & 0 & 1 & 2 & 3 & 4 & ... & 25\\
    \hline
    Caractère & A & B & C & D & E & ... & Z\\
    \hline
\end{tabular}

\par Lorsqu'un caractère est lu, son indice est émis, puis ce caractère est placé en première position et tous les autres caractères décalés.

\subsubsection{Algorithme}
\begin{minted}{c++}
/*@function MTF.
 *@param string chaine, la chaine à encoder.
 *@param vector<char> Alphabet, l'alphabet.
 *@return vector<size_t> indices, la chaine encodé dans une suite d'indices.
 */
vector<size_t> MTF(string chaine, vector<char> Alphabet)
{
	vector<size_t> indices;
	size_t k(0);

	for(string::iterator it = chaine.begin(); it != chaine.end(); it++)
	{
		for(size_t i(0); i < Alphabet.size(); ++i)
		{
			if(*it == Alphabet[i])
			{
				indices.push_back( i );
				char c = Alphabet[i];
				for( int z = i - 1; z >= 0; z-- )
	    			Alphabet[z + 1] = Alphabet[z];
 
        			Alphabet[0] = c;
			}
		}
	}
	
	return indices;
}
\end{minted}


\subsubsection{Spécificité}
\par Lorsque des caractères semblables se suivent, le flux émis contiendra beaucoup de 0 alors que dans une compression statistique (type codage de Huffman) augmentera considérablement le Taux de compression de données.

\subsubsection{Exemple :}
Exemple de l'application MTF à yabadabadoo\$ :

\begin{tabular}{|c|c|c|}
    \hline
    Caractère lu & Indices & Alphabet \\
    \hline
      &  & \$abcdefghijklmnopqrstuvwxyz \\
      \hline
    y & 25 & y\$abcdefghijklmnopqrstuvwxz\\
    \hline
     ya & 25;2 & ay\$bcdefghijklmnopqrstuvwxz\\
    \hline
     yab & 25;2;3 & bay\$cdefghijklmnopqrstuvwxz\\
    \hline
     yaba & 25;2;3;1 & aby\$cdefghijklmnopqrstuvwxz\\
    \hline
     yabad & 25;2;3;1;5 & daby\$cefghijklmnopqrstuvwxz \\
    \hline
     yabada & 25;2;3;1;5;1 & adby\$cefghijklmnopqrstuvwxz \\
    \hline
     yabadab & 25;2;3;1;5;1;2 & bady\$cefghijklmnopqrstuvwxz\\
    \hline
     yabadaba & 25;2;3;1;5;1;2;1 & abdy\$cefghijklmnopqrstuvwxz \\
    \hline
     yabadabad & 25;2;3;1;5;1;2;1;2 & daby\$cefghijklmnopqrstuvwxz \\
    \hline
     yabadabado & 25;2;3;1;5;1;2;1;2;16 & odaby\$cefghijklmnpqrstuvwxz \\
    \hline
     yabadabadoo & 25;2;3;1;5;1;2;1;2;16;0 & odaby\$cefghijklmnpqrstuvwxz \\
    \hline
     yabadabadoo\$ & 25;2;3;1;5;1;2;1;2;16;0;5 & \$odabycefghijklmnpqrstuvwxz \\
    \hline
\end{tabular}
        \subsection{La transformée de Burrows-Wheeler}

\par L'algorithme BWT, pour \textit{Burrows-Wheeler Transform} est un prétraitement utilisé en compression de données. Il ne s'agit pas d'un algorithme de compression, car aucune réduction de taille n'est effectuée. 
\par Il s'agit d'une méthode de réorganisation des données : le probabilité que des caractères identiques initialement éloignés les uns des autres se retrouves côte à côte dans le résultat est accrue.


\subsubsection{Fonctionnement}
\par Dans un premier temps on va copier la chaîne de caractère à coder dans un tableau carré en décalant la chaîne d'un caractère vers la droite à chaque nouvelle ligne. Ensuite, on va trier ces lignes par ordre alphabétique.
\par Grâce à ce décalage, les dernière lettre de chaque ligne précède la première lettre de la même ligne, sauf pour la ligne originale dont on notera la position. De plus comme les ligne sont rangées par ordre alphabétique, on peut retrouver la première colonne du tableau grâce à la dernière colonne.\\

\subsubsection{Algorithme du codage}
\begin{minted}{c++}
pair<size_t, string> BTW_encode(string chaine)
{
	
	string s;
	pair<int, string> result;
	vector<string> matrix;
	
	matrix.push_back(chaine);
	s.insert(0, chaine);
	
	for(size_t i(1); i < chaine.length(); ++i)
	{
		
		s.insert(s.begin(), s.end()-1, s.end());
		s.erase(s.end()-1, s.end());
		matrix.push_back(s);
		
	}
	
	sort(matrix.begin(), matrix.end());
	
	for(size_t i(0); i < matrix.size(); ++i)
	{
		
		result.second.insert(result.second.begin()+i, 
		                  matrix[i].end()-1, matrix[i].end());
		if(matrix[i][0] == 'y')
			result.first = i;

	}
	
	return result;	
}
\end{minted}


\subsubsection{Spécificité}
On observe que la transformée de textes relativement longs comportent de nombreuses répétitions contiguës de caractères car ils contiennent un nombre élevé de fois les mêmes mots. Cela provient du fait que la transformée de Burrows et Wheeler revient à trier les caractères par ordre alphabétique, puis à chaque fois écrire le caractère qui précédait dans le texte original. Vu qu'elle est construite par rotation, la dernière colonne est composée des caractères précédant. Un exemple flagrant est la transformation de « TEXTUELTEXTUEL » qui donne : « 7UUTTEELLXXTTEE ».

À des fins de compression donc, on utilise un algorithme de type MTF sur cette transformée. Les répétitions de caractères produiront beaucoup de zéros. Un codage de Huffman sur un tel résultat permet d'obtenir un taux élevé de compression. 

\subsubsection{Exemple :}
\par Prenons l'exemple, \textit{yabadabadoo\$} : \\ 

\paragraph{On fait la rotation :}

\begin{tabular}{|c|c|}
    \hline
    Indice & Chaine \\
    \hline
    0 & yabadabadoo\$\\
    \hline
    1 & \$yabadabadoo\\
    \hline
    2 & o\$yabadabado\\
    \hline
    3 & oo\$yabadabad\\
    \hline
    4 & doo\$yabadaba\\
    \hline
    5 & adoo\$yabadab\\
    \hline
    6 & badoo\$yabada\\
    \hline
    7 & abadoo\$yabad\\
    \hline
    8 & dabadoo\$yaba\\
    \hline
    9 & adabadoo\$yab\\
    \hline
    10 & badabadoo\$ya\\
    \hline
    11 & abadabadoo\$y\\
    \hline
\end{tabular}


\paragraph{On trie les chaînes :}


\begin{tabular}{|c|c|}
    \hline
    Indice & Chaine \\
    \hline
    0 & \$yabadabadoo\\
    \hline
    1 & abadabadoo\$y  \\
    \hline
    2 & abadoo\$yabad\\
    \hline
    3 & adabadoo\$yab\\
    \hline
    4 & adoo\$yabadab\\
    \hline
    5 & badabadoo\$ya\\
    \hline
    6 & badoo\$yabada\\
    \hline
    7 & dabadoo\$yaba\\
    \hline
    8 & doo\$yabadaba\\
    \hline
    9 & o\$yabadabado\\
    \hline
    10 & oo\$yabadabad\\
    \hline
    11 & yabadabadoo\$\\
    \hline
\end{tabular}

\paragraph{On repère l'indice de la chaîne à encoder et on l'encode}

\par Chaîne composée des derniers caractères des chaînes triées : oydbbaaaaod\$.
La chaînes \textit{yabadabadoo\$} apparaît à l'indice 11. On a donc le couple 11 et \textit{oydbbaaaaod\$} ou  simple la chaîne \textit{11oydbbaaaaod\$}.

\subsubsection{Décodage Burrows-Wheeler}

\par Pour décoder un code généré par BWT, il suffit de reconstruire le tableau des chaînes triées à l'aide de la chaîne fourni dans le code.

\par Pour reconstruire ce tableau il suffit de concaténer à l'avant des chaînes du tableau les caractères du code et de trier à chaque tour de boucle (nombre de caractères dans le code). On va donc avoir initialement dans notre tableau un caractères par chaînes (triées) fournis par le code. 

\paragraph{Avec \textit{oydbbaaaaod\$}}

\begin{tabular}{|c|c|}
    \hline
    Indice & Chaine \\
    \hline
    0 & o\\
    \hline
    1 & y  \\
    \hline
    2 & d\\
    \hline
    3 & b\\
    \hline
    4 & b\\
    \hline
    5 & a\\
    \hline
    6 & a\\
    \hline
    7 & a\\
    \hline
    8 & a\\
    \hline
    9 & o\\
    \hline
    10 & d\\
    \hline
    11 & \$\\
    \hline
\end{tabular}

\paragraph{Avec \textit{oydbbaaaaod\$} triée}

\begin{tabular}{|c|c|}
    \hline
    Indice & Chaine \\
    \hline
    0 & \$\\
    \hline
    1 & a  \\
    \hline
    2 & a\\
    \hline
    3 & a\\
    \hline
    4 & a\\
    \hline
    5 & b\\
    \hline
    6 & b\\
    \hline
    7 & d\\
    \hline
    8 & d\\
    \hline
    9 & o\\
    \hline
    10 & o\\
    \hline
    11 & y\\
    \hline
\end{tabular}

\paragraph{On concatène chaque caractère de \textit{oydbbaaaaod\$} à nos chaînes}

\begin{tabular}{|c|c|}
    \hline
    Indice & Chaine \\
    \hline
    0 & o\$\\
    \hline
    1 & ya  \\
    \hline
    2 & da\\
    \hline
    3 & ba\\
    \hline
    4 & ba\\
    \hline
    5 & ab\\
    \hline
    6 & ab\\
    \hline
    7 & ad\\
    \hline
    8 & ad\\
    \hline
    9 & oo\\
    \hline
    10 & do\\
    \hline
    11 & \$y\\
    \hline
\end{tabular}

\paragraph{On trie à nouveau}

\begin{tabular}{|c|c|}
    \hline
    Indice & Chaine \\
    \hline
    0 & \$y\\
    \hline
    1 & ab   \\
    \hline
    2 & ab \\
    \hline
    3 & ad\\
    \hline
    4 & ad\\
    \hline
    5 & ba\\
    \hline
    6 & ba\\
    \hline
    7 & da\\
    \hline
    8 & do\\
    \hline
    9 & o\$\\
    \hline
    10 & oo\\
    \hline
    11 & ya\\
    \hline
\end{tabular}

\paragraph{et ainsi de de suite}

\begin{tabular}{|c|c|}
    \hline
    Indice & Chaine \\
    \hline
    0 & o\$y\\
    \hline
    1 & yab   \\
    \hline
    2 & dab \\
    \hline
    3 & bad\\
    \hline
    4 & bad\\
    \hline
    5 & aba\\
    \hline
    6 & aba\\
    \hline
    7 & ada\\
    \hline
    8 & ado\\
    \hline
    9 & oo\$\\
    \hline
    10 & doo\\
    \hline
    11 & \$ya\\
    \hline
\end{tabular}

\begin{tabular}{|c|c|}
    \hline
    Indice & Chaine \\
    \hline
    0 & \$ya\\
    \hline
    1 & aba   \\
    \hline
    2 & aba \\
    \hline
    3 & ada\\
    \hline
    4 & ado\\
    \hline
    5 & bad\\
    \hline
    6 & bad\\
    \hline
    7 & dab\\
    \hline
    8 & doo\\
    \hline
    9 & o\$y\\
    \hline
    10 & oo\$\\
    \hline
    11 & yab\\
    \hline
\end{tabular}

\paragraph{Une fois la table complète}

\par Il suffit de regarder la ligne qui correspond à l'indice donné par le code.


\subsubsection{Algorithme du décodage}
\begin{minted}{c++}
string BTW_decode(pair<size_t, string> code)
{
	
	vector<string> matrix;
	size_t k(0);
	
	for(auto e : code.second)
		matrix.push_back(string(1, e));
		
	sort(matrix.begin(), matrix.end());
	
	for(size_t i(1); i < code.second.length(); ++i)
	{
		
		k = 0;
		
		for(auto &e : matrix)
		{
			
			e.insert(e.begin(), code.second[k]);
			++k;
			
		}

		sort(matrix.begin(), matrix.end());
		
	}

	return matrix[code.first];
	
}
\end{minted}
        \subsection{Codage arithmétique}

\par Le codage arithmétique est un codage \textbf{entropique} de compression de données sans perte. Il permet une meilleure compression que le codage de \textbf{Huffman}, sauf lorsque tous les poids pour les feuilles/nœuds/racines de l'arbre de \textbf{Huffman} sont des puissances de 2, auquel cas les deux méthodes sont équivalentes.

\subsubsection{Fonctionnement}
\par Dans un premier temps on va calculer la probabilité d'apparition pour chaque symbole. A partir de ces probabilités on va subdiviser l'intervalle $[0,1[$ pour chaque symbole en fonction de leurs probabilités d'apparition. Par exemple si le symbole \textbf{e} apparaît à 30\% dans notre texte, alors son intervalle fera 0.3.
\par A partir de cet intervalle et des probabilités d'apparitions on va pouvoir encoder.
\\

\subsubsection{Algorithme du codage}

\begin{algorithm}[H]
\caption{Codage arithmétique}
Pour chaque symbole, l'associer avec sa fréquence d'apparition.\;
BI = 0\; 
BS = 1\; 
\While{Il reste un caractère c du texte à coder}
{
	\For{s symbole de S}
	{    Attribuer un intervalle à s dans [BI, BS]
	}
	   
	BI = BIc\;
	BS = BSc\;
}
Choisir un n entre BI et BS\;
Renvoyer n\;
\end{algorithm}


\subsubsection{Spécificité}
Ce qui différencie le codage arithmétique des autres codages sources est qu'il code le message par morceaux (théoriquement il peut coder un message entier de taille quelconque mais dans la pratique on ne peut coder que des morceaux d'une quinzaine de symboles en moyenne) et représente chacun de ces morceaux par un nombre n (flottant) là où \textbf{Huffman} code chaque symbole par un code précis. Le problème qui en résulte pour le codage \textbf{Huffman} est qu'un caractère ayant une probabilité très forte d'apparition sera codé sur au moins un bit. Par exemple, si on cherche à coder un caractère représenté à 90\%, la taille optimale du code du caractère sera de 0,15 bit alors que \textbf{Huffman} codera ce symbole sur au moins 1 bit, soit 6 fois trop3. C'est cette lacune que le codage arithmétique comble grâce à un algorithme proche du \textbf{codage par intervalle}. 

\subsubsection{Exemple :}
\par Prenons l'exemple, \textit{canadian} : \\ 

\begin{tabular}{|c|c|}
    \hline
    Symbole & Probabilités \\
    \hline
    a & 0.375\\
    \hline
    c & 0.125\\
    \hline
    d & 0.125\\
    \hline
    i & 0.125\\
    \hline
    n & 0.25\\
    \hline

\end{tabular}

\par On encode à partir du tableau précédant :

\begin{tabular}{|c|c|}
    \hline
    Symbole & Intervalle \\
    \hline
     & [0, 1[\\
    \hline
    c & [0.375, 0.5[\\
    \hline
    a & [0.375, 0.421875[\\
    \hline
    n & [0.41015625, 0.421875[\\
    \hline
    a & [0.41015625, 0.41455078125[\\
    \hline
    d & [0.412353515625, 0.41290283203125[\\
    \hline
    i & etc\\
    \hline
\end{tabular}

Le mot \textbf{canadian} peut s'éncoder par \textbf{0.41271972656}

\subsubsection{Décodage arithmétique}

\par Pour décoder un code généré par codage arithmétique il suffit d'avoir la table des symboles avec leurs fréquences d'apparitions ainsi que la taille du texte. 

\par A partir de ces informations on va donc pouvoir reconstruire un "découpage" en intervalle dans [0, 1[ pour chaque symbole. A chaque tour $n = (n - BIc) / p$ où \textbf{n} est le réel, \textbf{BIc} la borne inférieur du caractère trouvé et \textbf{p} la probabilité d'apparition de ce même caractère.  

\subsubsection{Exemple : }

\par Prenons l'exemple suivant \textbf{0110100110101} avec comme taille 8 et la probabilité par symbole suivante :

\begin{tabular}{|c|c|}
    \hline
    Symbole & Probabilités \\
    \hline
    a & 0.375\\
    \hline
    c & 0.125\\
    \hline
    d & 0.125\\
    \hline
    i & 0.125\\
    \hline
    n & 0.25\\
    \hline
\end{tabular}

On décode \textbf{0110100110101} ce qui donne : 
 $0*2^{-1}+1*2^{-2}+1*2^{-3}+0*2^{-4}+1*2^{-5}+0*2^{-6}+0*2^{-7}+1*2^{-8}+1*2^{-9}+0*2^{-10}+1*2^{-11}+0*2^{-12}+1*2^{-13} = 0.41271972656$
\\
\begin{tabular}{|c|c|c|c|}
    \hline
    Réel & intervalle & symbole & position\\
    \hline
    0.41271972656 & [0.375, 0.5[ & c & 0\\
    \hline
    0.3017 & [0.0, 0.375[ & a & 1\\
    \hline
    0.8046 & [0.75, 0.1[ & n & 2\\
    \hline
    0.2186 & [0.0, 0.375[ & a & 3\\
    \hline
    0.5831 & [0.5, 0.625[ & d & 4\\
    \hline
    0.6653 & [0.625, 0.75[ & i & 5\\
    \hline
    0.3227 & [0.0, 0.375[ & a & 6\\
    \hline
    0.8606 & [0.75, 0.1[ & n & 7\\
    \hline
   
\end{tabular}

On retrouve comme mot : canadian.

\section{Théorie algorithmique de l'information}
	
	La théorie algorithmique de l'information, initiée par Kolmogorov, Solomonov et Chaitin dans les années 1960, vise à quantifier et qualifier le contenu en information d'un ensemble de données, en utilisant la théorie de la calculabilité et la notion de machine universelle de Turing.

	\subsection{Mesures de l'information}

\begin{definition}{Machine}
%comp. = composition ?
\par Composition de fonctions partielles prenant une chaîne et renvoyant une chaîne.
\end{definition}

\begin{definition}{M-Description}
\par Soit une machine $M$ telle que $M(\sigma) = x$. $\sigma$ est une $M-Description$ de $x$.
\end{definition}

\paragraph{Chaîne binaire et nombre entier} Les fonctions utilisées sont de $\mathbb{N}$ dans $\mathbb{N}$, mais on représentera les nombres entiers sous forme de chaînes binaires. Ainsi, cela devient des fonctions de $\{0,1\}^*$ dans $\{0,1\}^*$.
\par Ainsi, $\sigma$ sera la représentation sous forme de chaîne binaire de l'entier naturel $n$. La représentation binaire de $n+1$ est $1\sigma$. $n = nombre(\sigma)$ et $\sigma = chaine(n)$.

\begin{definition}{$chaine(n)$}
\par $$ rec\_chaine(n) = 
			\left \{
				\begin{array}{ll}
					rec\_chaine(n/2).(n~mod~2) & si~ n > 1 \\
					1 & si~ n = 1 \\
					0 & sinon \\
				\end{array}
			\right.
	$$
\par $$ chaine(n) = 
			\left \{
				\begin{array}{ll}
					rec\_chaine(n-1) & si~ n > 0 \\
					\epsilon & sinon \\
				\end{array}
			\right.
	$$
\end{definition}

\begin{definition}{$nombre(\sigma)$}
\par $$ rec\_nombre(n) = 
			\left \{
				\begin{array}{ll}
					2*rec\_nombre(u) & si~ \sigma = u.0 \\
					2*rec\_nombre(u)+1 & si~ \sigma = u.1 \\
					0 &sinon \\
				\end{array}
			\right.
	$$
\par $$ nombre(n) = 
			\left \{
				\begin{array}{ll}
					rec\_nombre(\sigma)+1 & si~ \sigma \not= \epsilon \\
					0 & sinon \\
				\end{array}
			\right.
	$$
\end{definition}

\begin{example}
\begin{itemize}
\item La chaîne $10$ correspond au nombre $5$.
\item $\epsilon$ correspond à $0$.
\end{itemize}
\end{example}

\begin{definition}{Complexité}
\par La complexité est la taille de la plus petite chaîne $\sigma$ permettant d'engendrer la chaîne $x$ pour une machine $M$.
\par $$C_M(x) = min\{|\sigma|:M(\sigma) = x\}$$
\par $min \emptyset = \inf$.
\end{definition}

	\subsection{Jetons de Rosza}
		\subsubsection{Fonctions récursives primitives}

\begin{definition}
En théorie de la calculabilité, une fonction récursive primitive est une fonction construite à partir de la fonction nulle, fonction successeur, fonctions projections et les schémas primitifs de récursion et de composition. Ces fonctions constituent un sous-ensemble strict des fonctions récursives.
\par Elles ont été initialement analysées par la mathématicienne Rózsa Péter.
\par On construit les fonctions récursives primitives par induction en partant des trois fonctions de base :
\begin{itemize}
	\item La fonction identiquement nulle ;
	\item La fonction successeur : $Succ(t)=t+1$ ;
	\item Les projections : $(x_{1},...,x_{k})\rightarrow x_{i}$.
\end{itemize}
et en itérant les deux constructions suivantes :
\begin{itemize}
	\item La composition de fonctions : 
si $g_{1}, g_{2}, ..., g_{k}$ sont récursives primitives sur $\mathbb {N} ^{n}$ 
et si $h$ est récursive primitive sur $\mathbb {N} ^{k}$, toutes déjà définies,
alors la fonction $f=h(g_{1},...,g_{k})$ est une fonction récursive primitive définie sur $\mathbb {N} ^{n}$ ;

	\item La définition récursive d'une fonction : 
si $g$ est récursive primitive sur $\mathbb {N} ^{n}$, et $h$ récursive primitive sur ${\mathbb  N}\times {\mathbb  N}\times {\mathbb  N}^{n}$, on définit une nouvelle fonction récursive primitive $f$ sur ${\mathbb  N}^{{n+1}}$ par :
	\begin{itemize}
		\item $f(0,\vec y) = g(\vec y)$
		\item $ f(Succ(x),\vec y) = h(x,f(x,\vec y),\vec y)$.
	\end{itemize}
\end{itemize}
\end{definition}

\begin{definition}{Jetons de bases}
\begin{itemize}
\item \Zero : $~ \mapsto 0$
\item \Successor : $~x \mapsto x+1$
\item \Identity : $~x \mapsto x$
\end{itemize}
\end{definition}

\begin{definition}{Constructeurs}
\begin{itemize}
\item \Left : $~x, \bar{y} \mapsto \bar{y}$
\item \Right : $~\bar{x}, y \mapsto \bar{x}$
\item Soient une fonction $f$ récursive primitive d'arité $k$ et les fonctions récursives primitives $g_1,g_2,\ldots,g_k$ d'arité $n$. \par $\Composition(f,g_1,g_2,\ldots,g_k)(\bar{x}) = f(g_1(\bar{x}),g_2(\bar{x}),\ldots,g_k(\bar{x}))$.
\item Soit $f$ une fonction récursive primitive d'arité $n$ et $g$ une fonction récursive primite d'arité $n+2$. On obtient, la fonction récursive primitive d'arité $n+1$ :
\par $\Recurse(f,g)(x, \bar{y}) = \left\{
	\begin{array}{ll}
		f(\bar{y}) & si ~x = 0 \\
		g(x-1,\Recurse(f,g)(x-1, \bar{y}),\bar{y}) & sinon
	\end{array} \right.
	$
\end{itemize}
\end{definition}

	\subsubsection{Fonctions $\mu$-recursives}
	
\begin{definition}
En informatique et en mathématiques, le terme fonction récursive (ou fonction $\mu$-recursive) désigne une classe de fonctions calculables, autrement dit de fonctions dont les valeurs peuvent être calculées à partir de leurs paramètres par un processus mécanique.
\par Les fonctions récursives primitives sont une sous-partie stricte de l'ensemble des fonctions récursives, elles ne permettent pas de calculer tout ce qui est calculable. Il est nécessaire d'y ajouter une fonction nommée $\mu$.
\end{definition}

\begin{definition}{Fonction $\mu$}
\par $\Mu(f) \bar{y} \mapsto min\{x : f(x,\bar{y}) = 0\}$.
\par À noter : contrairement aux fonctions récursives primitives, $\mu$ n'est pas définie sur tout son domaine : il s'agit d'une fonction récursive partielle.
\end{definition}

	\subsubsection{Interpréteur de jetons en ligne}

\href{http://rozsa-tokens.info/}{Interpréteur de jetons de Rozsa en ligne : http://rozsa-tokens.info/}

\subsection{Conditionnal information}

\begin{definition}
$$\mathbb{V}^2(0^{e-1}1\sigma,y)=\phi_e(<\sigma,y>)$$
\end{definition}

\begin{definition}
$$C(x|y) = min\{|\sigma|:\mathbb{V}^2(\sigma,y)=x\}$$
\end{definition}

\subsection{Lien entre C et K}

\par K est une complexité qui est plus naturelle (parce qu'on prend en compte qu'il y a de l'information dans la longueur).
\par Pour prouver $K(x) \leqslant^+ K(|x|)+|x|$, on a construit une machine sans préfixe N qui tente de décomposer l'entrée $\tau$ en $\tau=\sigma\rho$ tel que $\sigma = |\rho|$ : on va pouvoir en déduire un lien entre C et K.

\begin{theorem}
$$K(x)\leqslant K(C(x))+C(x)$$.
\end{theorem}

\begin{theorem}
$$K(X) \leqslant^+ K(|x|)+|x|$$
\end{theorem}

\begin{theorem}
$$C(X) \leqslant^+ K(x) \leqslant^+ C(x) + 2 log(C(x)) \leqslant^+ C(x)+2log(|x|)$$
\end{theorem}

\begin{proof}
On sait que $C(x) \leqslant^+ K(x).$
\end{proof}

\begin{theorem}
$$\forall x,y, C(xy) \leqslant^+ K(x) + C(y)$$
\end{theorem}

\begin{theorem}
$$abs(C(x)-C(y)) \leqslant^+ K(abs(x-y)) \leqslant^+ 2log(abs(x-y))$$
\end{theorem}

\begin{theorem}
$$C(x) =^+ K(x|C(x))$$
\end{theorem}

\begin{corollary}
$$C(xy) \leqslant^+ C(x) + 2 log(C(x))+C(y)$$
\end{corollary}
\section*{Annexe}
\subsection*{Sources}
\begin{itemize}
\item \href{https://www.sciencepresse.qc.ca/blogue/2014/12/24/mesurer-linformation-selon-hartley}{Mesurer l'information selon Hartley}
\item \href{https://fr.wikipedia.org/wiki/Entropie_de_Shannon}{Entropie de Shannon}
\item \href{https://centenaire-shannon.cnrs.fr/chapter/la-theorie-de-information}{Claude Shannon : Le monde en binaire}
\item \href{https://en.wikipedia.org/wiki/Kraft\%E2\%80\%93McMillan_inequality}{Kraft–McMillan inequality}
\item \href{https://fr.wikipedia.org/wiki/Codage_de_Shannon-Fano}{Codage de Shannon-Fano}
\item \href{https://en.wikipedia.org/wiki/Shannon\%E2\%80\%93Fano_coding}{Shannon–Fano coding}
\item \href{https://fr.wikipedia.org/wiki/Th%C3%A9orie_algorithmique_de_l%27information}{Théorie algorithmique de l'information}
\item \href{https://fr.wikipedia.org/wiki/Complexit%C3%A9_de_Kolmogorov}{Complexité de Kolmogorov}
\item \href{https://fr.wikipedia.org/wiki/Fonction_r%C3%A9cursive_primitive}{Fonction récursive primitive}
\item \href{https://fr.wikipedia.org/wiki/Fonction_r%C3%A9cursive}{Fonction récursive}
\item \href{https://en.wikipedia.org/wiki/%CE%9C-recursive\_function}{$\mu$-recursive function}
\item \href{https://fr.wikipedia.org/wiki/Move-to-front}{Move-to-front}
\item \href{https://fr.wikipedia.org/wiki/Transform\%C3\%A9e_de_Burrows-Wheeler}{Transformée de Burrows-Wheeler}
\end{itemize}


\end{document}